% =============================================================================
% Chapter 16: Discussion -- Part III (Extended Validation)
% =============================================================================
\mychapter{16. Discussion --- Extended Validation}

The single most important result in this part is the generalisation gap itself
(Section~15.2), not any one preprocessing or masking choice. Validation Dice in the 0.60--0.73 range
was already established as a plausible, defensible detection result by Part II; this part shows that
number does not predict performance on a sample genuinely never seen during training, which instead
falls to 0.12--0.21 across every configuration tested, including the original Part II configuration
itself (0.600 vs.\ 0.147, Section~10.3). Because this pattern reproduces across two independently run
experiments with different sample-role assignments, different loss functions, and different
preprocessing pipelines, it is best explained as a property of this pipeline's train/validation/test
design, specifically, how few distinct samples are available to fill the validation and held-out-test
roles, rather than as an artefact of any single training run. This reframes the acceptance-criteria
shortfall already stated in Part II (Section~10.2) as a structural finding rather than a training
issue that more epochs or a different loss function would resolve on its own.

Within that constraint, the preprocessing and masking ablations (Sections~15.2, 15.3) still show a
genuine, monotonic improvement in the more trustworthy of the two numbers: held-out test Dice rises
from 0.1218 to 0.1693 to 0.2119 as BHC/ring correction and then shape-aware masking are added, even as
validation Dice stays flat or drifts slightly downward. Read together with the generalisation-gap
finding above, this supports two separate conclusions rather than one: the pipeline's absolute
detection accuracy on unseen data remains far below the acceptance criteria, and the specific
preprocessing and masking changes tested here are moving that number in the right direction rather
than the wrong one, evidence that the improvements are real generalisation gains and not artefacts of
fitting the validation sample more closely.

The loss-function ablation (Section~15.1) resolves an inconsistency in how Part II's own headline
result and its stated primary loss function relate to one another: the held-out generalisation-gap
checkpoint's own recorded validation Dice (0.7275, epoch 13) is an exact match to the loss-ablation's
BCE row, meaning the thesis's central generalisation finding was already produced using BCE, not
combined Dice-Focal, the loss formulation Part II states as primary. Combined with BCE's higher raw
Dice and IoU and a recall difference against Dice-Focal (0.0072) too small to support the earlier
precision/recall trade-off argument for preferring Dice-Focal, this is read here as sufficient grounds
to recommend BCE as the pipeline's primary loss function, correcting rather than merely supplementing
Part II's stated choice.

The conv-block attribution result (Section~15.6) offers one concrete, if bounded, architectural
implication: with encoder\_1--3 together responsible for roughly two-thirds of the network's
attributed defect-discriminating signal and the final decoder block the least influential of all nine,
the highest-leverage architectural change for this task is more likely to be strengthening the early
encoder stages or their skip connections, for instance via attention gating, than adding depth to the
bottleneck or decoder. This is offered as a direction for future architecture work rather than a
finding validated by ablation; a single attribution method on a single representative crop is
suggestive evidence of where signal concentrates, not proof that modifying those specific blocks would
improve held-out performance.

The boundary-detection work (Sections~15.3--15.5) illustrates a pattern already established once in
this thesis, in Section~9.5/10.5's diagnosis of Bernsen's DCT-calibration failure, and repeated here in
a different subsystem: a method that produces a superficially plausible number can still be wrong by a
large factor, and the wrongness is only caught by an independent cross-check. sample\_06's 2.92\%
uncorrected defect fraction was a physically plausible-looking value for a metal AM part; only the
erosion-margin sweep revealed it was a boundary artefact rather than genuine porosity. Canny edge
detection (Section~15.5) supplies the same lesson from the opposite direction: it independently
confirms the circular-fit method's numbers on samples where both methods work, which is useful
corroboration, but it was not adopted for sample\_06/07 despite otherwise working correctly through
most of each volume's depth, because a method that is right most of the time but silently wrong in a
way that looks correct is a worse production choice than a method validated not to exhibit that
failure mode at all.

Finally, the two findings of Sections~14.9--14.10 and~14.10 respectively strengthen and complicate the
thesis's use of the NIST CoCr dataset. The gravimetric cross-check in Kim et al.'s own paper means the
$r = 0.995$ external-validation figure (Table~\ref{tab:mask_r_comparison}) is agreement with a
reference that already carries independent physical grounding, a stronger claim than XCT-vs-XCT
agreement alone would support, even with the caveat that individual per-sample confirmation of that
check was not verified here. Conversely, the voxel-size finding, samples 3 and 5 scanned at roughly a
third of the other samples' voxel size, raises a plausible, previously unconsidered explanation for why
those specific two samples were the ones excluded from this thesis's own training and ranking
(Sections~9.4, 9.5): a substantially finer scan resolution implies a smaller physical field of view per
slice at the same pixel dimensions, which could independently affect both defect content captured and
intensity/contrast statistics. This is stated explicitly as an unconfirmed hypothesis rather than a
demonstrated cause, consistent with this thesis's practice elsewhere of separating what has been
measured from what remains plausible but unverified, but it is a concrete, checkable direction for
follow-up work that this part's own investigation did not have scope to pursue further.
